%%%%%%%%%%  scan idx 629  |  folio 616  %%%%%%%%%%
\SGAPrintPage{616}%
\sgaStableTarget{sga6.e13.expose}\section*{\underline{EXPOS\'E XIII}\\[3pt] \underline{\underline{FINITENESS THEOREMS FOR THE PICARD FUNCTOR}}}
\begin{center}
\underline{by S. KLEIMAN}
\end{center}

\sgaStableTarget{sga6.e13.sec.1}\subsection*{1. The (b)-sheaves}
Let $k$ be an algebraically closed field and $X$ a projective $k$-scheme, endowed with an ample invertible module ${O}_X(1)$.

\sgaStableTarget{sga6.e13.def.1-1}\textbf{Definition 1.1 (Mumford).} \emph{Let $m\in\mathbb Z$. A coherent sheaf $F$ on $X$ is said to be \emph{$m$-regular} if, at every point of the support of $F$, ${O}_X(1)$ is generated by its global sections and if, for every $q\ge1$, we have}
\[
 H^q(F(m-q))=0.
\]

\sgaStableTarget{sga6.e13.lem.1-2}\textbf{Lemma 1.2.} \emph{Let \(0\longrightarrow F(-1)\longrightarrow F\longrightarrow G\longrightarrow 0\) be an exact sequence of coherent sheaves on $X$}\SGAAuthorFootnote{*}{Here, as below in analogous situations, it is understood that \(F(-1)\to F\) is defined by tensoring with a given section of \(\underline{O}_X(1)\).}\emph{. If $F$ is $m$-regular, then so is $G$.}

Indeed, we have the corresponding exact sequence:
\[
H^q(F(m-q))\longrightarrow H^q(G(m-q))\longrightarrow H^{q+1}(F(m-q-1)).
\]

\sgaStableTarget{sga6.e13.prop.1-3}\textbf{Proposition 1.3.} \emph{If a sheaf $F$ on $X$ is $m$-regular, then, for every $n\ge m$, we have}~:
\begin{enumerate}[label=(\roman*)]
\item \emph{$F$ is $n$-regular.}
\item \emph{\(H^0(F(n))\otimes H^0({O}_X(1))\longrightarrow H^0(F(n+1))\) is surjective.}
\item \emph{$F(n)$ is generated by \(H^0(F(n))\).}
\end{enumerate}

%%%%%%%%%%  scan idx 630  |  folio 617  %%%%%%%%%%
\SGAPrintPage{617}%
The proof is by induction on the dimension \(s\) of the support of \(F\); when \(s=-1\), it is trivial. Take a section \(\sigma\in H^0({O}_X(1))\) which does not vanish at any point of the finite set \(\Ass(F)\); then \(\sigma\) defines an exact sequence \(0\longrightarrow F(-1)\longrightarrow F\longrightarrow G\longrightarrow 0\), and \(G\) is \(m\)-regular by 1.2., with support of dimension \(s-1\).

Consider the corresponding exact sequence:
\[
H^q(F(n-q-1))\longrightarrow H^q(F(n-q))\longrightarrow H^q(G(n-q)).
\]
By the induction hypothesis, we have \(H^q(G(n-q))=0\) for \(n\ge m\), so \(H^q(F(n-q-1))=0\) implies \(H^q(F(n-q))=0\); but \(H^q(F(m-q))=0\), which gives (i).

To establish (ii), we chase through the following commutative diagram:
\[
\begin{tikzcd}[column sep=1.8em,row sep=2.5em]
 & & & H^0(F(n))\otimes H^0({O}_X(1)) \arrow[rr] \arrow[d]
 & & H^0(G(n))\otimes H^0({O}_X(1)) \arrow[r] \arrow[d]
 & 0 \\
0 \arrow[rr]
& & H^0(F(n)) \arrow[ur] \arrow[r]
& H^0(F(n+1)) \arrow[rr]
& & H^0(G(n+1)) \arrow[d] \\
& & & & & 0
\end{tikzcd}
\]
where the oblique arrow is defined by \(\varphi\mapsto\varphi\otimes\sigma\), the horizontal rows are exact, the first by the \(n\)-regularity of \(F\) and by (i), already proved, and the second column is exact by the induction hypothesis.

Finally, (iii) follows from the following commutative diagram:
\[
\begin{tikzcd}[column sep=large,row sep=large]
H^0(F(n))_X\otimes H^0({O}_X(1))_X
  \arrow[r,"\alpha_n"]
  \arrow[d]
& H^0(F(n)(1))_X
  \arrow[d,"\beta_{n+1}"] \\
\bigl[H^0(F(n))_X\bigr](1)
  \arrow[r,"\beta_n\otimes\mathrm{id}"]
& F(n)(1)
\end{tikzcd}
\]
%%%%%%%%%%  scan idx 631  |  folio 618  %%%%%%%%%%
\SGAPrintPage{618}%
where, if \(M\) is a vector space over \(k\), we write \(M_X\) for \(M\otimes_k{O}_X\). For \(n\ge m\), \(\alpha_n\) is surjective by (ii); therefore, if \(\beta_{n+1}\) is surjective, so is \(\beta_n\). By Serre's theorem, \(\beta_{n+1}\) is surjective for \(n\gg0\); hence (iii).

\sgaStableTarget{sga6.e13.prop.1-4}\textbf{Proposition 1.4.} \emph{Let \(0\to F(-1)\to F\to G\to0\) be an exact sequence on \(X\). If \(G\) is \(m\)-regular, then we have}~:
\begin{enumerate}[label=(\roman*)]
\item \emph{\(H^q(F(n))=0\) for \(q\ge2\) and \(n\ge m-q\).}
\item \emph{\(h^1(F(n-1))\ge h^1(F(n))\) for \(n\ge m-1\).}
\item \emph{\(H^1(F(n))=0\) for \(n\ge (m-1)+h^1(F(m-1))\).}
\end{enumerate}
\emph{In particular, \(F\) is \([m+h^1(F(m-1))]\)-regular.}

Indeed, in the exact sequence
\[
H^{q-1}(G(n))\longrightarrow H^q(F(n-1))\longrightarrow H^q(F(n))\longrightarrow H^q(G(n)),
\]
the two outer groups are \(0\) for \(q\ge2\) and \(n\ge m-(q-1)\) by 1.3 (i). Thus \(H^q(F(m-q))\xrightarrow{\sim}H^q(F(m-q+1))\xrightarrow{\sim}H^q(F(m-q+2))\xrightarrow{\sim}\ldots\), but, by Serre's theorem, \(H^q(F(n))=0\) for \(n\gg0\); hence (i).

For \(n\ge m-1\), consider the following exact sequence:
\[
0\to H^0(F(n-1))\to H^0(F(n))
  \xrightarrow{\alpha_n}H^0(G(n))\to H^1(F(n-1))
  \to H^1(F(n))\to0
\]
It implies (ii). Now suppose that \(\alpha_n\) is surjective and consider the following commutative diagram:
\[
\begin{tikzcd}[column sep=large,row sep=large]
H^0(F(n))\otimes H^0({O}_X(1))
  \arrow[r,"\alpha_n\otimes\mathrm{id}"]
  \arrow[d]
& H^0(G(n))\otimes H^0({O}_X(1))
  \arrow[r]
  \arrow[d,"\beta_n"]
& 0 \\
H^0(F(n+1))
  \arrow[r,"\alpha_{n+1}"]
& H^0(G(n+1)).
\end{tikzcd}
\]
%%%%%%%%%%  scan idx 632  |  folio 619  %%%%%%%%%%
\SGAPrintPage{619}%
If \(n\ge m\), then \(\beta_n\) is surjective by 1.3 (ii); therefore \(\alpha_{n+1}\) is surjective. Consequently, \(H^1(F(n-1))\xrightarrow{\sim}H^1(F(n))\xrightarrow{\sim}\ldots\xrightarrow{\sim}0\). On the other hand, if \(\alpha_n\) is not surjective, then \(h^1(F(n-1))>h^1(F(n))\); therefore \(n\mapsto h^1(F(n))\) is \uline{strictly} decreasing from \(h^1(F(m-1))\) until it becomes \(0\); hence (iii).

\sgaStableTarget{sga6.e13.def.1-5}\textbf{Definition 1.5.} \emph{Let \(F\) be a coherent sheaf on \(X\), \(r\) an integer \(\ge\operatorname{dim{.}supp}(F)\) and \((b)=(b_0,\ldots,b_r)\) a sequence of \(r+1\) integers. Then \(F\) is said to be a \emph{\((b)\)-sheaf} if it satisfies one of the following two equivalent sets of conditions}~:

\sgaStableTarget{sga6.e13.para.1-5-1}\((1.5.1)\)
\begin{enumerate}[label=(\roman*)]
\item \emph{At every point of the support of \(F\), \({O}_X(1)\) is generated by \(H^0({O}_X(1))\).}
\item \emph{\(h^0(F(-1))\le b_0\).}
\item \emph{(When \(r\ge1\)), there exists a section \(\sigma\in H^0({O}_X(1))\) defining an exact sequence \(0\to F(-1)\to F\to G\to0\) such that \(G\) is a \((b_1,\ldots,b_r)\)-sheaf.}
\end{enumerate}

\sgaStableTarget{sga6.e13.para.1-5-2}\((1.5.2)\) \emph{There exists an \(F\)-regular sequence of \(r\) sections \(\sigma_1,\ldots,\sigma_r\in H^0({O}_X(1))\) such that, for \(i=0,\ldots,r\), if \(F_i\) denotes the restriction of \(F\) to the zero scheme of \(\sigma_1,\ldots,\sigma_i\), then \(h^0(F_i(-1))\le b_i\).}

\sgaStableTarget{sga6.e13.prop.1-6}\textbf{Proposition 1.6.} \emph{Let \(F\) be a \((b)\)-sheaf on \(X\). Then}~:
\begin{enumerate}[label=(\roman*)]
\item \emph{For every ``sufficiently general'' sequence \(\sigma\) of \(r\) sections \(\sigma_1,\ldots,\sigma_r\in H^0({O}_X(1))\), if \(F_{\sigma,i}\) denotes the restriction of \(F\) to the zero scheme of \(\sigma_1,\ldots,\sigma_i\), then \(h^0(F_{\sigma,i}(-1))\le b_i\), where \(0\le i\le r\).}
\item \emph{Every subsheaf \(G\) of \(F\) is a \((b)\)-sheaf.}
\end{enumerate}

%%%%%%%%%%  scan idx 633  |  folio 620  %%%%%%%%%%
\SGAPrintPage{620}%
Indeed, let \(S'\) be the affine space \(\mathbb A_k^N\) whose \(k\)-points correspond to the sequences \(\sigma\), and \(T\) the open subset of \(S'\) corresponding to the \(F\)-regular sequences. For fixed \(i\), the sheaves \(F_{\sigma,i}(-1)\) corresponding to the \(k\)-points of \(T\) are contained in a flat family over \(T\), and by hypothesis \(T\) is nonempty. Thus (i) follows from the upper semi-continuity property of the function $\sigma\longmapsto h^0(F_{\sigma,i}(-1))$ (EGA III, 7.7.5). Finally, since every ``sufficiently general'' sequence is \(G\)-regular and such that \(G_{\sigma,i}\) is a subsheaf of \(F_{\sigma,i}\) for every \(i\), (ii) follows from (i).

\sgaStableTarget{sga6.e13.lem.1-7}\textbf{Lemma 1.7.} \emph{Let \(0\to F(-1)\to F\to G\to0\) be an exact sequence on \(X\). If the Hilbert polynomial of \(F\) is written in the form}
\[
 \chi(F(n))=\sum_{i=0}^{r}a_i\binom{n+i}{i},
\]
\emph{then}
\[
 \chi(G(n))=\sum_{i=0}^{r-1}a_{i+1}\binom{n+i}{i}.
\]
Indeed, \(\chi(G(n))=\chi(F(n))-\chi(F(n-1))=\sum_{i=0}^{r}a_i\left[\binom{n+i}{i}-\binom{n-1+i}{i}\right]\); therefore \(\chi(G(n))=\sum_{i=1}^{r}a_i\binom{n+i-1}{i-1}\), which gives the assertion.

\sgaStableTarget{sga6.e13.prop.1-8}\textbf{Proposition 1.8.} \emph{Let \(F\) be a \((b)\)-sheaf on \(X\), \(s=\operatorname{dim{.}supp}(F)\) and \(\chi(F(n))=\sum_{i=0}^{s}a_i\binom{n+i}{i}\) the Hilbert polynomial of \(F\). Then}~:
\begin{enumerate}[label=(\roman*)]
\item \emph{For \(n\ge -1\), \(h^0(F(n))\le \sum_{i=0}^{s}b_i\binom{n+i}{i}\).}
\item \emph{\(a_s\le b_s\) and \(F\) is also a \((b_0,\ldots,b_{s-1},a_s)\)-sheaf.}
\end{enumerate}
Indeed, we argue by induction on \(s\). When \(s=0\), this is clear since \(a_0=h^0(F)=h^0(F(-1))\le b_0\). When \(s\ge1\), we have an exact sequence
\[
0\longrightarrow F(-1)\longrightarrow F\longrightarrow G\longrightarrow0
\]
